% Original: PDF strana 27, donji slajd.
\slajdid{02-s054}
\begin{frame}[label=02-s054]{Hazardi i usklađivanje kašnjenja}
\setlength{\parskip}{0pt}
% element: 02-s054:e01
\Oznaka{Proizvod zbirova}
Isti postupak važi kada minimizaciju radimo prekrivanjem \textbf{logičkih nula}. Pri prelazu između susednih površina koje obezbeđuju nulu može se pojaviti \alert{\textbf{lažna jedinica}}. Dodajemo površinu koja ih povezuje i ima zajednička polja sa obe.\par
\par\vspace{3mm}
% element: 02-s054:e02
\Oznaka{Kada se menja više ulaza}
U praksi je problem složeniji: može se istovremeno menjati \textbf{više ulaznih signala}. Oni stižu u I/ILI deo kombinacione mreže sa \textbf{različitim kašnjenjima}. Povezivanje površina za promenu jednog ulaza zato ne rešava sve takve prelaze.\par
\par\vspace{3mm}
% element: 02-s054:e03
\Oznaka{Usklađivanje kašnjenja i preuređivanje mreže}
Kašnjenja ujednačavamo dodavanjem komponenti, npr. invertora, ili preuređujemo mrežu i \alert{\textbf{odstupamo od minimalne realizacije}}. To se često radi heuristički ili simulacijom. Presudni su \textbf{razumevanje i iskustvo projektanta}.
\end{frame}
